% ============================================================
%  chapters/app-a-constraint-spaces.tex
%  Appendix A: Sets, Relations, and Constraint Spaces
% ============================================================

\section*{A.1 Introduction}

Many structures developed throughout this monograph share a common pattern.
One begins with a large space of possibilities, imposes constraints, identifies
equivalent configurations, and studies the resulting admissible structures.
Whether one is discussing logical operators, gesture recognition, keyboard
motion, semantic inference, cosmological trajectories, or nested Spherepop
expressions, the same mathematical architecture repeatedly appears.

The purpose of this appendix is to formalize that architecture.

\begin{tcolorbox}[enhanced, colback=RSVPdeep, colframe=RSVPdeep,
  coltext=white, arc=6pt, boxrule=0pt,
  left=6mm, right=6mm, top=5mm, bottom=5mm]
  \textbf{Central thesis.} Structure is not fundamental.
  Constraint is fundamental. Objects emerge as stable equivalence classes
  within constrained spaces.
\end{tcolorbox}

% ── A.2 Constraint Spaces ─────────────────────────────────────
\section*{A.2 Constraint Spaces}

\begin{definition}{Constraint Space}{constraint-space-app}
  A \emph{constraint space} is a triple $\mathcal{C} = (X, A, \kappa)$
  consisting of:
  \begin{itemize}
    \item a \emph{state space} $X$,
    \item an \emph{admissible subset} $A \subseteq X$,
    \item a \emph{constraint functional}
      $\kappa : X \to \mathbb{R}_{\geq 0}$
  \end{itemize}
  such that $A = \kappa^{-1}(0)$. States with $\kappa(x) = 0$ satisfy all
  constraints; states with $\kappa(x) > 0$ violate one or more constraints.
\end{definition}

\begin{example}
  \textbf{Arithmetic expressions.} Let
  $X = \{ 1{+}2,\ (1{+}2),\ (3),\ 3,\ ((3)) \}$.
  Define $\kappa(x)$ as the number of unevaluated arithmetic operations.
  Then $\kappa(1{+}2) = 1$ while $\kappa(3) = 0$.
  The admissible set consists precisely of fully reduced expressions.
\end{example}

\begin{example}
  \textbf{Unit circle.} Let $X = \mathbb{R}^2$ and
  $\kappa(x,y) = |x^2 + y^2 - 1|$.
  Then $A = \{(x,y) \in \mathbb{R}^2 : x^2 + y^2 = 1\}$.
\end{example}

% ── A.3 Constraint Descent ────────────────────────────────────
\section*{A.3 Constraint Descent}

\begin{definition}{Constraint Descent Flow}{constraint-descent}
  A \emph{constraint descent flow} is a dynamical system
  \[
    \frac{dx}{dt} = -\nabla\kappa(x).
  \]
\end{definition}

\begin{proposition}{Monotone Descent}{monotone-descent}
  Along a constraint descent flow, $\kappa(x(t))$ is nonincreasing.
\end{proposition}

\begin{proof}
  Differentiate:
  $\frac{d}{dt}\kappa(x(t)) = \nabla\kappa \cdot \frac{dx}{dt}$.
  Substituting $\frac{dx}{dt} = -\nabla\kappa$ gives
  $\frac{d}{dt}\kappa = -\|\nabla\kappa\|^2 \leq 0$.
\end{proof}

\begin{conceptaside}{Constraint Descent as a Unifying Principle}
  This result appears throughout the monograph in many guises:
  logical reduction, Spherepop evaluation, semantic refinement,
  RSVP relaxation, optimization, and learning. All are instances of
  constraint descent in appropriately defined spaces.
\end{conceptaside}

% ── A.4–A.5 Equivalence and Quotients ─────────────────────────
\section*{A.4 Equivalence Relations}

\begin{definition}{Equivalence Relation}{equiv-rel}
  An \emph{equivalence relation} $\sim$ on $X$ satisfies reflexivity
  ($x \sim x$), symmetry ($x \sim y \Rightarrow y \sim x$), and
  transitivity ($x \sim y$, $y \sim z \Rightarrow x \sim z$).
\end{definition}

\begin{example}
  Arithmetic evaluation induces $((3)) \sim (3) \sim 3$.
  All denote the same value.
\end{example}

\begin{example}
  In gesture recognition, two hand trajectories may differ microscopically
  while producing the same letter: $g_1 \sim g_2$ whenever they correspond
  to the same recognized symbol.
\end{example}

\section*{A.5 Quotient Spaces}

\begin{definition}{Quotient Space}{quotient-space-app}
  The \emph{quotient space} is
  $X/{\sim} = \{[x] : x \in X\}$
  where $[x] = \{y : y \sim x\}$.
\end{definition}

This is the mathematical foundation of symbolic compression, gesture
canonicalization, semantic normalization, and Spherepop collapse operators.

\begin{proposition}{Quotient Projection is Surjective}{quot-surj}
  The quotient projection $\pi : X \to X/{\sim}$ defined by $\pi(x) = [x]$
  is surjective.
\end{proposition}

\begin{proof}
  For any $[a] \in X/{\sim}$, choosing $x = a$ gives $\pi(a) = [a]$.
\end{proof}

% ── A.7 Central Theorem ───────────────────────────────────────
\section*{A.7 Objects as Equivalence Classes}

\begin{theorem}{Symbolic Representation Theorem}{symbolic-rep}
  Every symbolic system can be represented as an equivalence class within
  a richer state space.
\end{theorem}

\begin{proof}
  Let $S$ be a symbolic system. Construct a state space $X$ consisting of
  all possible realizations of symbols. Define $x \sim y$ iff $x$ and $y$
  represent the same symbol. Then equivalence classes $[x]$ correspond
  exactly to symbolic objects, giving $S \cong X/{\sim}$.
\end{proof}

\begin{remark}
  The letter ``A'' is not the ink. It is the equivalence class of all
  acceptable realizations. A gesture is not a trajectory. It is the
  equivalence class of trajectories recognized as equivalent.
\end{remark}

% ── A.8 Central Principle ─────────────────────────────────────
\section*{A.8 The Constraint–Quotient Principle}

\begin{namedthm}{Constraint–Quotient Principle}
  Structure emerges from the interaction of (1) a possibility space $X$,
  (2) constraints $\kappa$, (3) equivalence relations $\sim$, and
  (4) projections $\pi$. The resulting object is $(X, \kappa, \sim, \pi)$.
  Everything in this monograph may be viewed as a particular realization
  of this architecture:
  \[
    \text{Constraint}
    \;\to\;
    \text{Equivalence}
    \;\to\;
    \text{Projection}
    \;\to\;
    \text{Structure.}
  \]
\end{namedthm}

\medskip
\noindent
\textit{Appendix~B} will show how projection necessarily creates information
loss. \textit{Appendix~C} will show that Spherepop evaluation is a
constraint-descent process on expression trees.
